\section{Ordered-word coefficient families}

\begin{theorem}[Polynomial coefficients and word frequencies]
    \label{thm:word_coefficient_expansion}
    \lean{Polynomial.coeff_pow_sum_smul_add_C_mul_X_eq_sum_counts}
    \leanok
    Let $\mathcal A$ be a unital complex algebra, let $I$ be a finite alphabet,
    and let $H_i,T\in\mathcal A$. Adjoin a distinguished letter $\ast$ to $I$,
    with associated operator $V_\ast=T$ and $V_i=H_i$ for $i\in I$.
    Fix integers $m,k\ge0$. For a word $w$ of length $m+k$, write $n_i(w)$
    for the number of occurrences of $i$ and put
    \begin{align*}
        O(w)=V_{w(0)}V_{w(1)}\cdots V_{w(m+k-1)}.
    \end{align*}
    For each $a\in\mathbb N^I$ with $\sum_i a_i=m$, define
    \begin{align*}
        f_a=\sum_{\substack{w:\ n_\ast(w)=k\\n_i(w)=a_i\ (i\in I)}}O(w).
    \end{align*}
    If $X$ is a central indeterminate, then for every $z\in\mathbb C^I$,
    \begin{align*}
        [X^k]\left(\sum_{i\in I}z_iH_i+TX\right)^{m+k}
            =\sum_{\substack{a\in\mathbb N^I\\\sum_i a_i=m}}
                \left(\prod_{i\in I}z_i^{a_i}\right)f_a.
    \end{align*}
    The formula includes an empty alphabet and zero powers; the empty word
    contributes the identity. This identifies the distinguished-letter family
    in the proof of~\cite[Lemma~4.2]{Arad2013AreaLaw}.
\end{theorem}
\begin{proof}
    \leanok
    \uses{thm:finite_sum_power_ordered_words,def:poisson_word_count_vector}
    Expand the power into ordered words. A word contributes
    $X^{n_\ast(w)}$ times its unchanged ordered operator product, multiplied by
    $\prod_i z_i^{n_i(w)}$. Taking the coefficient of $X^k$ retains precisely
    the words with $n_\ast(w)=k$. Their ordinary-letter counts sum to $m$.
    Partition this finite sum by those counts and factor out the common scalar
    on each part. This gives the displayed identity, with every word retained
    with its original multiplicity.
\end{proof}

\begin{definition}[Full-frequency word coefficients]
    \label{def:full_word_coefficient}
    \lean{WordMultiplicity.wordCoefficient,WordMultiplicity.countsFinsupp_equiv}
    \leanok
    Let $I$ be finite and let $H_a$ belong to a semiring. For $d\ge0$ and
    $\beta:I\to\mathbb N$, define
    \begin{align*}
        C_{d,\beta}=\sum_{\substack{w\in I^{\{0,\ldots,d-1\}} \\n_w=\beta}}
        H_{w(0)}\cdots H_{w(d-1)}.
    \end{align*}
    A bijection $e:I\to J$ preserves letter counts:
    $n_{e\circ w}(e(a))=n_w(a)$.
    These are the coefficients in \cite[Lemma~4.2]{Arad2013AreaLaw}.
\end{definition}

\begin{definition}[Insertion of a distinguished frequency]
    \label{def:insert_marker_frequency}
    \lean{WordMultiplicity.insertMarkerFrequency,WordMultiplicity.sum_insertMarkerFrequency}
    \leanok
    Fix $u\in I$, put $O=I\setminus\{u\}$, and let $k\ge0$.
    For $\alpha:O\to\mathbb N$, define $\iota_{u,k}(\alpha):I\to\mathbb N$
    by its value $k$ at $u$ and its restriction $\alpha$ to $O$.
    Its total frequency is $k+\sum_{a\in O}\alpha(a)$.
\end{definition}

\begin{theorem}[A full-frequency slice in the residual alphabet]
    \label{thm:one_marker_full_frequency_identity}
    \lean{WordMultiplicity.wordCoefficient_insertMarkerFrequency,
        WordMultiplicity.sum_smul_wordCoefficient_marker_eq}
    \leanok
    \uses{def:full_word_coefficient,def:insert_marker_frequency,def:one_marker_word_coefficient}
    Let the operators $H_a$ belong to a unital complex algebra and let $m,k\ge0$.
    Put $E_m^O=\{\alpha:O\to\mathbb N:\sum_a\alpha(a)=m\}$, and let
    $D_{m,k,\alpha}$ be the ordered-word coefficient on the ordinary alphabet
    $O$, with $H_u$ occurring $k$ times.
    Then $C_{m+k,\iota_{u,k}(\alpha)}=D_{m,k,\alpha}$, and every complex weight
    function on the full frequency vectors satisfies
    \begin{align*}
        \sum_{\substack{\beta:I\to\mathbb N \\\sum_a\beta(a)=m+k,\ \beta(u)=k}}
        \lambda_\beta C_{m+k,\beta}
        =\sum_{\alpha\in E_m^O}
        \lambda_{\iota_{u,k}(\alpha)}D_{m,k,\alpha}.
    \end{align*}
    This identifies the actual coefficient family used in
    \cite[Lemma~4.2]{Arad2013AreaLaw}.
\end{theorem}
\begin{proof}
    \leanok
    \uses{def:full_word_coefficient,def:insert_marker_frequency}
    The bijection $O\sqcup\{*\}\to I$ sends $*$ to $u$ and fixes every
    ordinary letter. It gives a bijection of the words in the two coefficient
    sums, preserving their ordered products and all occurrence counts.
    Insertion of the distinguished count is a bijection from $E_m^O$ to the
    full-frequency slice. Summing the coefficient identity along this bijection
    gives the weighted identity.
\end{proof}
